% Generated draft: requires mathematical and source review.
\documentclass[11pt]{article}
\usepackage[margin=25mm]{geometry}
\usepackage{fontspec}
\setmainfont{Latin Modern Roman}
\usepackage{amsmath,amssymb,mathtools,unicode-math}
\setmathfont{Latin Modern Math}
\usepackage{enumitem,needspace,xurl,hyperref}
\setlength{\parindent}{0pt}
\setlength{\parskip}{5pt}
\setlength{\emergencystretch}{3em}
\newcommand{\N}{\mathbb{N}}
\newcommand{\Z}{\mathbb{Z}}
\newcommand{\Q}{\mathbb{Q}}
\newcommand{\R}{\mathbb{R}}
\newcommand{\C}{\mathbb{C}}
\newcommand{\F}{\mathbb{F}}
\newcommand{\A}{\mathbb{A}}
\newcommand{\PP}{\mathbb{P}}
\newcommand{\E}{\mathbb{E}}
\newcommand{\eps}{\varepsilon}
\newcommand{\dd}{\,\mathrm{d}}
\begin{document}
\section*{UnsolvedMath ID: 30001750}
\textbf{Source problem code:} catalogue entry
\subsection{Definitions and mathematical statement}
The supplied record concerns bandlimited signals \(f\) and \(g\), sampling expansions with coefficients \(\beta_k\), and samples \(f(k\varepsilon)\). Bandlimitedness ordinarily means bounded support of a Fourier transform; the source extract does not specify the signal domain, Fourier-transform convention, function space, or bandwidth. It also supplies neither the sampling expansions nor the relationship between \(f\), \(g\), and the coefficients.

For notation only, let \(I\) denote the unspecified coefficient and sample index set, and write
\[
\beta=(\beta_k)_{k\in I},\qquad S_{\varepsilon}f=(f(k\varepsilon))_{k\in I},
\]
whenever these expressions are defined. Here \(\varepsilon\) denotes the sampling step appearing in the supplied statement; its admissible values are unspecified. These expressions introduce notation without selecting an admissible signal class or asserting that the coefficients are determined by the samples.

\textbf{Recorded recovery questions.}
\begin{enumerate}
\item For the intended sampling expansions, can \(\beta\) be recovered stably, or approximated stably, from \(S_{\varepsilon}f\), and can this improve reconstruction?
\item Under suitable additional signal assumptions, can sampling below the Nyquist rate still permit recovery of \(\beta\), \(f\), and \(g\)?
\end{enumerate}
The signal assumptions, expansions, coefficient space, sampling space, stability norms or topologies, approximation criterion, reconstruction criterion, and bandwidth-dependent Nyquist threshold are absent. Consequently, these questions cannot be promoted from this extract to a fully quantified mathematical conjecture or theorem.

\textbf{Unresolved source qualification.} The supplied dated literature assessment says that subsampling has a positive answer under extra signal hypotheses while stable coefficient recovery and approximation remain a further question. A separate supplied statement-verification field instead says that the report answers both questions affirmatively under explicit hypotheses and that the extraction is not an open problem. Those hypotheses and answers are not reproduced. The registry labels the record \texttt{partially\_solved} and its statement \texttt{unrecoverable}; these are supplied curation labels, not independently verified mathematical conclusions.
\subsection{Short English statement}
The record asks whether coefficients in unspecified bandlimited sampling expansions can be recovered or approximated stably from samples, and whether additional assumptions allow recovery of the coefficients and both signals below the Nyquist rate. Essential definitions and hypotheses are missing, and the supplied descriptions disagree about which questions the report resolves. This draft preserves the questions for review without asserting an open conjecture.
\subsection{Sources}
\begin{itemize}
\item Supplied local UnsolvedMath registry record, permanent ID \(30001750\), title \emph{Stable Coefficient Recovery from Bandlimited Samples}, source problem code \texttt{catalogue entry}, in the Oberwolfach Reports Open Problems collection. The proposed year is recorded as 2011; no individual proposer is supplied.
\item Supplied report citation: \emph{Operator Algebras and Representation Theory: Frames, Wavelets and Fractals} (2011), Oberwolfach Reports, pp.~901--978, DOI 10.4171/owr/2011/17. Source URL and supplied literature-source link labelled ``original OWR report'': \url{https://doi.org/10.4171/owr/2011/17}. This authoritative report reference is supplied, but its contents have not been independently consulted for this draft.
\item Literature assessment dated exactly \texttt{2026-08-21}: supplied dated curation, not independently verified. It records a positive subsampling answer under extra signal hypotheses and leaves stable coefficient recovery and approximation as a further question. The supplied background requests independent verification of time-sensitive status and cited claims.
\item The separate supplied statement-verification note says the report answers both recovery questions affirmatively under explicit hypotheses, so the extraction is not an open problem; it retains the original record for provenance without promoting a separate clean statement. This conflicts with the scope of the dated assessment above and requires source review.
\end{itemize}
\end{document}
